% ============================================================
% Section 10: 19th Century Foundations (Slides 71-78)
% ============================================================
\section{19th Century Foundations}

% ----------------------------------------------------------
% Slide 71: Section Divider
% ----------------------------------------------------------
{
\setbeamercolor{background canvas}{bg=modernSteel}
\begin{frame}[plain]
  \centering
  \vspace{2cm}
  \textcolor{white}{\Huge\bfseries The Century of Rigor}\\[0.6em]
  \textcolor{white}{\Huge\bfseries and Revolution}\\[1em]
  \textcolor{white!80}{\large 1800--1900 CE}\\[1.5em]
  \textcolor{white!60}{\normalsize Galois \(\cdot\) Riemann \(\cdot\) Cantor \(\cdot\) Noether \(\cdot\) Ramanujan \(\cdot\) Kovalevskaya}

  \note{
    Transition from Analysis: ``Gauss bridges two centuries. What followed was a
    crisis: mathematicians asked what mathematics really IS.''
    The 19th century is where mathematics became rigorous --- and where some of its
    most romantic and tragic stories took place.
    Gauss (1777--1855) was covered in the previous section but his influence extends
    deeply into this era.
  }
\end{frame}
}

% ----------------------------------------------------------
% Slide 72: Evariste Galois -- Genius Cut Short
% ----------------------------------------------------------
\begin{frame}{Evariste Galois --- Genius Cut Short}

  \begin{columns}[T]
    \begin{column}{0.52\textwidth}
      \begin{personbox}[Evariste Galois (1811--1832)]{modernSteel}
        \textbf{Origin:} Bourg-la-Reine, France\\[3pt]
        \textbf{Key work:} Founded \alert{group theory}. Proved there is no
        general algebraic formula for polynomials of degree~5 or higher,
        building on Niels Henrik Abel (1802--1829, Norway) who first proved
        the impossibility of solving the general quintic by radicals
        (Abel--Ruffini theorem, 1824).
        Galois went further: he characterised \emph{exactly which} equations
        can be solved by radicals, using the group structure of their roots.
      \end{personbox}

      \vspace{0.4em}
      \begin{insightbox}
        He died in a duel at age~20. The night before, he wrote down his
        mathematical ideas, scribbling: ``\emph{I have not time.}''
      \end{insightbox}
    \end{column}

    \begin{column}{0.45\textwidth}
      % TikZ: Simple group theory concept
      \visualplaceholder[5cm]{Symmetry group of a triangle}
    \end{column}
  \end{columns}

  \note{
    ``He died at 20 in a duel. The night before, he frantically wrote down his
    mathematical ideas, scribbling in the margins: `I have not time.' Those ideas
    revolutionised algebra.''
    Credit Abel explicitly: Abel proved quintic unsolvability slightly earlier
    (published 1824). The Abel--Ruffini theorem established the impossibility;
    Galois theory provided the deeper framework explaining why.
    \verifytag{Abel published 1824, refined 1826. Galois provided deeper framework.}

    PACING: Linger here. This is the emotional climax. Read the quote slowly.
    Let students feel the urgency of a 20-year-old writing mathematics he
    knows he will not live to publish.
  }
\end{frame}

% ----------------------------------------------------------
% Slide 73: Bernhard Riemann
% ----------------------------------------------------------
\begin{frame}{Bernhard Riemann --- Curved Spaces and a Million-Dollar Question}

  \begin{columns}[T]
    \begin{column}{0.5\textwidth}
      \begin{personbox}[Georg Friedrich Bernhard Riemann (1826--1866)]{modernSteel}
        \textbf{Origin:} Breselenz, Kingdom of Hanover, Germany;
        worked in G\"ottingen\\[3pt]
        \textbf{Key work:}
        \begin{itemize}\setlength\itemsep{2pt}
          \item \alert{Riemannian geometry} --- curved spaces of arbitrary
                dimension, essential for Einstein's general relativity
          \item The Riemann integral
          \item The \textbf{Riemann hypothesis} (1859) --- still unproven
                as of 2026, one of the Millennium Prize Problems
                worth \textbf{\$1\,million}
        \end{itemize}
      \end{personbox}
    \end{column}

    \begin{column}{0.47\textwidth}
      % Curved space visualization
      \visualplaceholder[5cm]{Flat grid}
    \end{column}
  \end{columns}

  \note{
    ``His hypothesis about prime numbers, stated in 1859, remains unsolved.
    There is a million-dollar prize waiting for whoever proves it.''
    Riemannian geometry became the mathematical language for Einstein's
    general relativity --- space itself is curved.
    Died young at 39 from tuberculosis.
  }
\end{frame}

% ----------------------------------------------------------
% Slide 74: Georg Cantor -- Infinity Has Sizes
% ----------------------------------------------------------
\begin{frame}{Georg Cantor --- Infinity Has Sizes}

  \begin{columns}[T]
    \begin{column}{0.48\textwidth}
      \begin{personbox}[Georg Ferdinand Ludwig Philipp Cantor (1845--1918)]{modernSteel}
        \textbf{Origin:} St.\ Petersburg, Russia (German--Danish parentage);
        worked in Halle, Germany\\[3pt]
        \textbf{Key work:} Founded \alert{set theory}. Proved that there
        are different sizes of infinity --- the real numbers are ``more
        infinite'' than the natural numbers. This was so controversial
        that leading mathematicians tried to block his career.
      \end{personbox}

      \vspace{0.4em}
      \begin{insightbox}
        He proved that infinity comes in sizes. Some infinities are bigger
        than others. This drove him to a breakdown --- and his critics
        nearly destroyed his career.
      \end{insightbox}
    \end{column}

    \begin{column}{0.49\textwidth}
      % TikZ: Cantor's diagonal argument
      \visualplaceholder[5cm]{Suppose we list all reals between 0 and 1}
    \end{column}
  \end{columns}

  \note{
    Walk through the diagonal argument step by step.
    ``He proved that infinity comes in sizes. Some infinities are bigger than others.
    This drove him to a breakdown --- and his critics nearly destroyed his career.''
    Leopold Kronecker was his main antagonist:
    ``God made the integers; all else is the work of man.''
  }
\end{frame}

% ----------------------------------------------------------
% Slide 75: Emmy Noether -- The Mother of Modern Algebra
% ----------------------------------------------------------
\begin{frame}{Emmy Noether --- The Mother of Modern Algebra}

  \begin{columns}[T]
    \begin{column}{0.52\textwidth}
      \begin{personbox}[Amalie Emmy Noether (1882--1935)]{modernSteel}
        \textbf{Origin:} Erlangen, Germany; fled to the US in 1933\\[3pt]
        \textbf{Key work:}
        \begin{itemize}\setlength\itemsep{2pt}
          \item \alert{Noether's theorem} (1918): symmetries in physics
                $\longleftrightarrow$ conservation laws
          \item Transformed abstract algebra --- ring theory, ideal theory
          \item Denied a proper faculty position for years because she was
                a woman
        \end{itemize}
      \end{personbox}

      \vspace{0.4em}
      \begin{insightbox}
        Einstein called her ``the most significant creative mathematical
        genius since higher education opened to women.''
      \end{insightbox}
    \end{column}

    \begin{column}{0.45\textwidth}
      % Noether's theorem visualized
      \visualplaceholder[5cm]{Noether's Theorem}

      \vspace{0.3em}
      \visualplaceholder[4cm]{Portrait of Emmy Noether\\(Wikimedia Commons)}
    \end{column}
  \end{columns}

  \note{
    ``She was denied a salary for years because of her gender. Her colleagues
    had to list her courses under a male professor's name. She is now
    considered one of the greatest algebraists in history.''
    Her theorem is one of the most profound results in mathematical physics ---
    every symmetry implies a conservation law.
    She fled Nazi Germany in 1933 and died in the US in 1935.

    WOMEN'S THREAD -- CONNECT THE PATTERN: ``Remember Hypatia? Murdered for
    her learning, 1,500 years ago. Kovalevskaya, two slides from now -- needed
    a sham marriage just to be allowed to study. And Noether, here, in the 20th
    century, denied a salary because of her gender. This is not ancient history.
    This is a pattern. And we will see, before this lecture ends, who finally
    broke through.'' [Forward-link to Mirzakhani, Slide 95.]
  }
\end{frame}

% ----------------------------------------------------------
% Slide 76: Srinivasa Ramanujan
% ----------------------------------------------------------
\begin{frame}{Srinivasa Ramanujan --- The Self-Taught Genius}

  \begin{columns}[T]
    \begin{column}{0.52\textwidth}
      \begin{personbox}[Srinivasa Ramanujan Aiyangar (1887--1920)]{modernSteel}
        \textbf{Origin:} Erode, Tamil Nadu, India; worked at Cambridge\\[3pt]
        \textbf{Key work:} Self-taught mathematical genius. Produced thousands
        of results in number theory, infinite series, and continued fractions,
        many without formal proof. His notebooks continue to yield new
        mathematics a century later. Collaborated with G.\,H.~Hardy
        at Cambridge.
      \end{personbox}

      \vspace{0.4em}
      % Taxicab number anecdote
      \begin{tcolorbox}[colback=modernSteel!5, colframe=modernSteel!50,
        fonttitle=\bfseries\small, title={The Taxicab Number}, rounded corners,
        boxrule=0.5pt, left=3pt, right=3pt, top=2pt, bottom=2pt]
        \small
        Hardy visited Ramanujan and mentioned his taxi was number 1729 ---
        ``a rather dull number.'' Ramanujan instantly replied:\\[2pt]
        ``\emph{No, it is a very interesting number; it is the smallest
        number expressible as the sum of two cubes in two different ways.}''\\[2pt]
        $1729 = 1^3 + 12^3 = 9^3 + 10^3$
      \end{tcolorbox}
    \end{column}

    \begin{column}{0.45\textwidth}
      \visualplaceholder[5cm]{Page from Ramanujan's notebooks\\(Cambridge University Library)}

      \vspace{0.5em}
      % Some Ramanujan formulas
      \visualplaceholder[5cm]{Diagram}
    \end{column}
  \end{columns}

  \note{
    ``With almost no formal training, he produced results that astonished the best
    mathematicians in the world. Hardy called their collaboration `the one romantic
    incident in my life.'\,''
    Ramanujan died at 32 of tuberculosis.
    His notebooks are still being studied --- some formulas were proved only decades
    later.

    PACING: The taxicab number story (1729) is the lightest moment in this
    heavy section. Let the students laugh. They need it.
  }
\end{frame}

% ----------------------------------------------------------
% Slide 77: Sofia Kovalevskaya
% ----------------------------------------------------------
\begin{frame}{Sofia Kovalevskaya --- First Woman Doctorate in Mathematics}

  \begin{columns}[T]
    \begin{column}{0.52\textwidth}
      \begin{personbox}[Sofia Vasilyevna Kovalevskaya (1850--1891)]{modernSteel}
        \textbf{Origin:} Moscow, Russia; worked in Stockholm, Sweden\\[3pt]
        \textbf{Key work:}
        \begin{itemize}\setlength\itemsep{2pt}
          \item First woman to obtain a doctorate in mathematics in modern
                Europe (G\"ottingen, 1874)
          \item Major contributions to analysis, partial differential
                equations, and mechanics
          \item The \alert{Kovalevskaya top} --- a new integrable case
                of rigid body rotation
          \item First woman appointed to a full professorship in northern
                Europe (Stockholm, 1889)
        \end{itemize}
      \end{personbox}

      \vspace{0.3em}
      \begin{insightbox}
        To get an education, she entered a marriage of convenience just
        to be allowed to travel abroad and study.
      \end{insightbox}
    \end{column}

    \begin{column}{0.45\textwidth}
      \visualplaceholder[4.5cm]{Portrait of Sofia Kovalevskaya\\(Wikimedia Commons)}

      \vspace{0.4em}
      % Kovalevskaya top diagram
      \visualplaceholder[5cm]{Three known integrable cases}
    \end{column}
  \end{columns}

  \note{
    ``To get an education, she had to enter a marriage of convenience just to be
    allowed to travel abroad and study. She became one of the greatest analysts
    of her century.''
    \verifytag{``First woman PhD in mathematics'' needs qualification as
    ``in modern Europe'' since specifics depend on definitions.}
    The Prix Bordin jury was so impressed they increased the prize money.
    She died of influenza at only 41.
  }
\end{frame}

% ----------------------------------------------------------
% Slide 78: 19th Century Overview
% ----------------------------------------------------------
\begin{frame}{The 19th Century --- An Overview}

  \begin{center}
    \visualplaceholder[5cm]{Five pillars}
  \end{center}

  \vspace{-0.3em}
  {\small\textcolor{modernSteel}{\textbf{Transition:}} \textit{``The 20th century
  would bring a crisis at the very foundations of mathematics\ldots''}}

  \note{
    Brief overview --- hit the five pillars quickly.
    Emphasise that ALL of these developments challenged the existing understanding
    of what mathematics IS.
    Transition: ``The 20th century opened with Hilbert's challenge: can we prove
    math is consistent? The answer was shocking.''
  }
\end{frame}
